\subsection{PART I - Adaptivity Analysis}
\label{sec:intro-methodology-adapt}

\subsubsection{Language Models}
\paragraph{Existing Methodology and Limitations}
The state-of-art data analysis algorithms are mostly implemented in Python equipped with the .panda library for database access, C with external database access library, etc.
These languages and external libraries are mainly focusing on efficiency and scalable computation over big data.
In terms of adaptivity, none of them taking care of this property w.r.t. these data analysis algorithms implemented.
However, the variety of the external API makes the formalization of data analysis algorithms and its adaptivity property over these languages difficult.

% \todo{Brief Describe Existing Language Models}
\paragraph{Our Method}
We identify some data analysis models which are usually used as the base of data analysis algorithms.
Based on these models, we design simple but scalable language based on standard while language that can be used to implement these models and formalize
the adaptivity in these algorithms, named {\tt QWhile}.
It is a standard while language enriched with the query request command and added labels on each program command.  


\subsubsection{Program Property Formalization}

\paragraph{Existing Methodology and Limitations}
There are many methodologies for formalizing a program's conceptual property. 
For example, the abstract interpretation based methodology which is originated from Cousot's influential paper~\cite{CousotH78, CousotC77} and intensively developed in~\cite{Mine06, ZanioliC11, MastroeniZ17, KincaidCBR18, Cousot19a, JanaHAGC20}, etc. This category of method does well in formalizing property from one perspective such as the dependency relation~\cite{Cousot19a}, program point reachability~\cite{MastroeniZ08}, etc.
The execution history and program simulation-based~\cite{Israel56, Kildall73}
are also widely used in defining
and formalizing the program's properties. While they have the same limitation as the abstract interpretation and it is not straightforward how to apply them to analyzing and formalizing the \emph{adaptivity}
for adaptive data analysis program as a \emph{quantitative-reachability} property from two perspectives.
There is also purely syntactic-based methodology originated from~\cite{Weiser84} and developed in~\cite{DenningD77} in defining a property (the data dependency relation), but it suffers the limitation of certain language and is not easy to generalize into other setting, especially the \emph{adaptivity}.
Because of its nature as a \emph{quantitative-reachability} property, defining \emph{adaptivity} requires us to defining both the dependency property and the quantitative property for the program.
And there are mainly three technique difficulties among existing formalization methodologies in defining the \emph{adaptivity} w.r.t. the two sub-properties.
\begin{enumerate}
\item
Existing definitions on the semantics concept of the data / value dependency relation don't
consider both control influence and value influence, such as~\cite{Cousot19a, DenningD77, AbadiBHR99, Mantel2004, CheneyAA11}.
However, both of these
two influences are considered as \emph{may-dependent} in the intuitive query dependency 
in adaptive data analysis programs.
\item In the state-of-art data dependency relation formalization, there isn't any technique taking into account
the dependency depth, specifically the quantitative information of the data dependency property. For example the works in~\cite{GoguenM84, AssafNSTT17, BartheDR11}.
However, this dependency depth is the key quantity in formalizing the intuitive \emph{adaptivity}.
\item The intuitive \emph{adaptivity} doesn't accumulate linearly through the dependency relation and 
dependency depth. There is only very limited works related to this non-linear quantity property analysis.
The non-monotonic resource cost analysis in~\cite{AlbertFR15} focus on analyzing the peak resource consumption of the heap structure, which is still different from the accumulation of the \emph{adaptivity}.
This requires us to design new formal model in order to define this property precisely.
% here isn't any research combining the two pieces of information. No need to mention the adaptivity analysis.
\end{enumerate}
%
\paragraph{New Methodology}
Our semantics-based formalization model defines this intuitive \emph{adaptivity} as a \emph{quantitative-reachability} program property semantically as follows.
\begin{enumerate}
   \item From the \emph{reachability} perspective, we first define the \emph{data dependency} relation based on our trace-based semantics.
  %  The semantics-based definition is inspired from~\cite{Cousot19a, ZanioliC11, CheneyAA11}.
   Our new definition capture both the control influence and value influence between data, which are never fully captured in previous dependency relation~\cite{Cousot19a, ZanioliC11, CheneyAA11}.
   It can be used to reason about the dependency between queries.

   \item Then from the \emph{quantitative} perspective, we define the \emph{dependency quantity} with respect to certain data
     to capture the dependency times between queries. We borrow the idea from the program complexity analysis areas~\cite{GulwaniJK09}, where the complexity bound can be repurposed to estimate the times that certain data is used.
   
     \item With the two property defined above, we construct a dependency graph by combining both the \emph{dependency relation} (as the graph edge) and \emph{dependency quantity} (as the vertex weight) on it. This dependency graph accurately reflects the two perspectives of the intuitive \emph{adaptivity}. Though there are previous works designing the weight dependency graph, none of them able to carry both dependency and dependency depth and apply it into analyzing the \emph{adaptivity}.
   
   \item Based on this graph, we formally define the \emph{adaptivity} as a \emph{quantitative-reachability} program property as the longest finite walk. This is also a novel quantity never considered before.
   \end{enumerate}
   % \item 
\subsubsection{Program Analysis for Adaptivity}
\label{sec:intro-static}

\paragraph{Existing Methodology and Limitations}
There are many static program analysis techniques, which are widely used in dependency, resource cost, quantitative analysis, etc. 
However, it is not straightforward how to apply them into analyzing and estimating the adaptivity -- an \emph{quantitative-reachability} property carrying two perspective sub-properties.
\begin{enumerate}
\item To the best of our knowledge,
existing static dependency analysis techniques do not consider both control influence and value influence.
However, to estimate the \emph{adaptivity} accurately, we need to consider both of them.
\item The existing static analysis techniques in the complexity or resource cost estimation areas,
cannot give accurate reachability times on programs 
at every execution location.
\item Moreover, existing program analysis focus
either only the dependency analysis
or only the program quantitative property (such as complexity or resource cost),
but not both of them.
\end{enumerate}

\paragraph{New Methodology}
We present a static program analysis framework, named {\THESYSTEM} in Chapter~\ref{sec:adapt-static} to estimate the adaptivity.
This new methodology reuses the techniques in data flow, control flow analysis, and reachability-bound analysis.
It gives tighter bounds on the adaptivity of a program than the ones one would achieve 
by directly using the data and control flow analyzes or the ones that one would achieve 
by directly using reachability bound analysis techniques alone. 
\begin{enumerate}
\item From the \emph{reachability} perspective, {\THESYSTEM} first 
% adopts both the static data flow analysis and the control flow analysis techniques for
estimates the \emph{data dependency} relation accurately. The estimation is more accurate than existing works in the sense that
it performs both the data flow analysis and the control flow analysis.

\item Then {\THESYSTEM} reasons the
\emph{quantitative} perspective through \emph{reachability-bound} analysis techniques.
It estimates the evaluation times w.r.t. each data. This estimation captures the \emph{quantitative} information for adaptivity, which is not able to be captured by any existing static program analysis framework.

\item
{\THESYSTEM} then constructs an \emph{estimated dependence} graph encoding both the estimated \emph{reachability} and \emph{quantitative} information.
This encoding complements the weakness of both lines of works in comparison to purely use single either \emph{reachability} or \emph{quantitative} analysis, 
% {\THESYSTEM} estimates both the two sub-properties of \emph{adaptivity} from \emph{reachability} and \emph{quantitative} perspectives more accurate.
% Moreover, the combination isn't trivial adoption and addition, we combine them effectively by encoding them on to a dependency graph and compute the \emph{adaptivity} over the graph in the next step.

%
\item In the last step, we design a new path search algorithm based on the graph theory 
and compute the \emph{adaptivity} accurately on this graph.
The computation fully utilize the two perspectives sub-properties estimated above.
This new algorithm again also complements the NP-hardness of the longest path search problem and computes an accurate symbolic length of the longest finite walk based on an efficient symbolic weighted longest path estimation.
\end{enumerate}

\subsection{PART II - The Path-sensitive Reachability-bound Analysis}
\label{sec:intro-methodology-reachability}

This property is a worst-case bound on the number of times a given control location 
inside a procedure is visited during the program execution.
Many works in the program analysis area develop algorithms estimating the program's loop bound or overall complexity.
However, there isn't a general analysis algorithm that solves this problem directly or path-sensitively.
To solve these issues and improve the analysis results on this problem, we present a new algorithm
aiming to find the precise symbolic \emph{reachability-bound} on the program's every control location.
Below are the methodologies we adopted in our new algorithm and their existing limitations.

\subsubsection{Program Abstraction}
\paragraph{Existing Techniques and Limitation.}
The program abstraction techniques are widely used in program analysis area to precisely summarize the key property of the program. They are usually based on some specific abstract interpretation models as introduced in Section~\ref{sec:intro-methodology-adapt}.
In analyzing the loop bounds, the unified lattice model~\cite{CousotH78}, polyhedra~\cite{CousotC77} and octagons~\cite{Mine06} are used in~\cite{GulwaniZ10} for summarize the loop transition and generating the transition system.
But none of the three abstraction models can represent the if guard in multipath loop, and solving the transition system using these abstraction models is low-efficient. The wedge abstract developed in~\cite{KincaidCBR18} for generating the non-linear loop invariant
computing the \emph{reachability-bound} can also be used. While it only works well for the specific targeting problem, and still low-efficient or inaccuracy in summarizing the multipath loops and solving the path-sensitive \emph{reachability-bound} problem. 

\paragraph*{Our Method.}
In our algorithm, we choose to use the difference constrain based program abstraction model combined with boolean expression abstraction to generate our abstract transition graph $\absG(c)$ for a program $c$. 
Each command in the program is abstracted into an abstract transition, $\absevent$.
It can represent the program loops path-sensitively which is more accurate than the state-of-art program abstraction models.

It is based on the difference constrain model in~\cite{SinnZV17} and the size change abstraction in~\cite{SinnZV14}. The efficiency also benefit from our
semantics model, which is inspired from~\cite{Cousot19a} and~\cite{Cousot19} and has similar intuition with the program model in~\cite{SinnZV17}.


\subsubsection{Amortized Complexity Analysis through Ranking Function}
\paragraph{Existing Techniques and Limitation.}
This method is originated from Tarjan's influential paper~\cite{PotechinP17} combined with \emph{ranking function} estimation from~\cite{BradleyMS05} and developed in~\cite{ZulegerGSV11, SinnZV14, SinnZV17, LuCT21, AliasDFG10}.
It does well in nested loops by alternating the loop bound computation with the ranking function estimation. While estimating the ranking function ignores the interleaving between multiple paths of one loop.
It over-approximates the bound of each single path as well the loop bound when the path interleaving affects the loop execution.
\paragraph{Our Method.}
We utilize the effectiveness of this method by first computing the ranking function $\locbound(\absevent, c)$ for each transition edge $\absevent$, and then estimating the bound on each ranking function, $\varinvar(\locbound(\absevent, c), c)$ and transition edge, $\absclr(\absevent, c)$ alternatively and path-insensitively.
This alternation lines up with advantage of the \emph{amortized complexity analysis} is efficient and accurate without recursively unrolling the nested loops when composing the bound of different paths.
But it is not precise enough without considering the effect of the path-interleaving.
In order to compute the \emph{reachability-bound} path sensitively, we also deploy the path refinement technique.

\subsubsection{Loop Summarization and Path Refinement}
\paragraph{Existing Techniques and Limitation.}
Another line of loop bound analysis through loop summarization and path refinement seek for precise loop path representation~\cite{ManoliosV06, BalakrishnanSIG09, SharmaDDA11, Flores-MontoyaH14, HumenbergerJK18, CyphertBKR19}, and compute loop bound over accurate loop summarization~\cite{GulwaniJK09, ZulegerGSV11}.
They do well in summarizing the loop path and computing the interleaving between paths, as well as computing the precise bound considering the interleaving. However, composing the bound between nested loops. Recursively unrolling the nested loops and computing the interleaving between unrolled paths are heavy and non-terminating.
\paragraph{Our Method.}
Our method take advantage from the accuracy of the path refinement techniques by a light-weight path refinement algorithm adapted from~\cite{GulwaniJK09}.
We compute the refined program $\rprog$ based on the abstract transition graph and then
effectively combine this methodology with the \emph{amortized complexity analysis} above through ranking function estimation.
% The next methodology introduced in the next part 

\subsubsection{Reachability-bound Computation}
\paragraph{Existing Techniques and Limitation.}
\cite{GulwaniZ10} first introduce \emph{reachability-bound} as the upper bound on the number of times a given control location 
inside a procedure is visited during program execution.
While existing works all have different limitations when inferring the \emph{reachability-bound}.
For example, Gulwani et. al~\cite{GulwaniZ10}
give a two-step solution by combining program abstract and proof-rules-based bound computation.
But it does not compute the \emph{reachability-bound} precisely for every control location.
Instead, they only compute the complexity bounds for the overall loop and over-approximate for different locations using this bound,
which is loose especially when there are multipath loops.
Techniques in program complexity analysis~\cite{GustafssonEL05,HumenbergerJK18} 
or worst-case resource cost analysis
~\cite{BrockschmidtEFFG16,AlbertAGP08,AliasDFG10,Flores-MontoyaH14} can be repurposed to estimate similar quantities such as the
bounds on loop iterations, execution time, etc.
While these quantities focus only on 
the overall complexity,
none of them compute the reachability-bound on a given program control location directly or path-sensitively.
This leads to inaccuracy in data leakage analysis, resource cost estimation, etc.

\paragraph{Our Method}
To overcome these limitations, 
our method is built over an abstract transition graph which effectively combines amortized complexity analysis with loop summarization-based multipath refinement.
This combination mitigates the limitations faced by either technique individually. 
Building on this combination, we introduce two novel quantities,
the \emph{path reachability-bound} and \emph{loop reachability-bound}.
They are the main novelty of our new path-sensitive reachability-bound algorithm build on the traditional methods.
It is computed over the refined program and effectively utilizing the ranking function estimated through the abstract transition graph.
These two bounds combines the two lines of techniques effectively and complements the weakness in both lines.
The first one bounds the evaluation times of each loop-free and interleave-free path in a refined program, and the second one bounds the number of iterations of an outer loop w.r.t. an inner loop, counting only the iterations when the inner loop is ``entered''. 
We present algorithms for estimating these two quantities. By using them we can then 
%Section~\ref{sec:looprb} and~\ref{sec:pathrb} 
compute the \emph{reachability-bound} for each program point in a path-sensitive way.
%in Section~\ref{sec:psrb}.
We have implemented our algorithm in a prototype and evaluated it on 290 program extracted from different benchmarks. Our evaluation %Section~\ref{sec:eval} 
shows that we can accurately estimate different bounds for different locations in programs with loops containing multiple paths. Based on this, we can also compute a tighter overall complexity compared to the state-of-art worst-case complexity analysis tools.
